\clearpage
\SourceLeaf{103}%

\begin{center}
\large [INVERSION DES ELLIPTISCHEN INTEGRALS ERSTER GATTUNG.]
\end{center}

\bigskip

\noindent\hspace*{\parindent}Über die transzendente Funktion $\displaystyle\int\frac{dx}{\BareRoot{(1-xx)(1-\mu xx)}}=\varphi$, wobei wir \mbox{$x=f\varphi$} setzen.

Setzen wir $\log x = y$, so haben wir

\[ \frac{x\,dy}{d\varphi} = \BareRoot{(1-xx)(1-\mu xx)}, \]

\[ \frac{ddy}{d\varphi^2} = -\frac{dx}{x\,d\varphi}\cdot\frac{dy}{d\varphi} - \frac{x\,dx}{(1-xx)\,d\varphi}\cdot\frac{dy}{d\varphi} - \frac{\mu x\,dx}{(1-\mu xx)\,d\varphi}\cdot\frac{dy}{d\varphi}, \]

\[ \frac{ddy}{d\varphi^2} = -\frac{(1-xx)(1-\mu xx)}{xx} - (1-\mu xx) - \mu(1-xx), \]

\[ \frac{ddy}{d\varphi^2} = -\frac{1}{xx} + \mu xx. \]

\noindent Es sei

\[ \int\frac{d\varphi}{xx} = t, \qquad \int \mu xx\,d\varphi = u, \]

\noindent dann ist

\[ \frac{dy}{d\varphi} = u - t. \]

\noindent Ferner sei

\[ \int u\,d\varphi = v, \qquad \int t\,d\varphi = w, \]

\noindent dann ist

\[ y = v - w. \]

\noindent Wir setzen

\[ e^{-w} = P\varphi, \qquad e^{-v} = Q\varphi, \]

\[ PP'' = P'P' - QQ, \qquad QQ'' = Q'Q' - \mu PP. \]

\[ P = \varphi - \tfrac{1}{6}(1+\mu)\varphi^3 + \tfrac{1}{120}(1+4\mu+\mu\mu)\varphi^5 - \cdots, \]

\[\begin{aligned} Q &= 1+\ast-\tfrac{1}{12}\mu\varphi^4+\tfrac{1}{90}(\mu+\mu\mu)\varphi^6\\&\quad-\tfrac{1}{10080}(8\mu+17\mu\mu+8\mu^3)\varphi^8-\cdots.\end{aligned}\]

\bigskip

